%% Based on the official SCP Original Software Publication template.
%% Copyright 2007, 2008, 2009 Elsevier Ltd.
%% This template derivative may be distributed under the LaTeX Project Public
%% License version 1.2 or later. See http://www.latex-project.org/lppl.txt .
%% The software and original prose have separate licensing; see repository notices.
\documentclass[preprint,12pt,a4paper]{elsarticle}
\usepackage{amssymb}
\usepackage{amsmath}
\usepackage{hyperref}
\hypersetup{hidelinks,pdftitle={OrchLang: Auditing Resource Contracts for Bounded LLM Workflows}}
\journal{Science of Computer Programming}
\newcommand{\evCppTests}{106}
\newcommand{\evPythonTests}{31}
\newcommand{\evCostWorkflows}{23}
\newcommand{\evCostRuns}{115}
\newcommand{\evGenerated}{576}
\newcommand{\evAccepted}{24}
\newcommand{\evAcceptedFiniteCounterexamples}{0}
\newcommand{\evFiniteEqualRejections}{16}
\newcommand{\evTraceIdentityDisagreements}{0}
\newcommand{\evMutantsDetected}{4}
\newcommand{\evTokenizerEncodings}{24}
\newcommand{\evUnderestimates}{13}
\newcommand{\evConcatenations}{10}
\newcommand{\evNonadditive}{6}
\newcommand{\evCertificates}{24}
\newcommand{\evLeanChecks}{624}
\newcommand{\evSourceTampers}{48}
\newcommand{\evMalformedBatches}{4}
\newcommand{\evReplayCases}{12}
\newcommand{\evReplayPaths}{48}
\newcommand{\evExternalProjections}{11}
\newcommand{\evExternalRuns}{220}
\newcommand{\evExternalBoundViolations}{0}
\newcommand{\evExternalBaselineDisagreements}{0}
\newcommand{\evOriginalApplicationExecutions}{0}
\newcommand{\evSpentUsd}{0.00}

\newcommand{\evDataflowSourcePrograms}{256}
\newcommand{\evDataflowAccepted}{16}
\newcommand{\evDataflowExecutionContexts}{3072}
\newcommand{\evDataflowAcceptedFiniteCounterexamples}{0}
\newcommand{\evDataflowSourceReconstructionDisagreements}{0}
\newcommand{\evDataflowLeanNormalizationDisagreements}{0}
\newcommand{\evDataflowLeanReferenceDisagreements}{0}
\newcommand{\evDataflowLeanReplayDisagreements}{0}
\newcommand{\evDataflowCppAccepted}{8}
\newcommand{\evDataflowCppSecretArgumentCases}{192}
\newcommand{\evDataflowBoundaryCases}{8}
\newcommand{\evDataflowBoundaryPaths}{23}
\newcommand{\evDataflowSourceTampersRejected}{32}
\newcommand{\evDataflowAlteredPlansRejected}{16}
\newcommand{\evDataflowInvalidEvaluatorInputsRejected}{8}
\newcommand{\evDataflowCppTests}{106}
\newcommand{\evDataflowPythonTests}{47}
\newcommand{\evDataflowProofFiles}{3}

\begin{document}
\begin{frontmatter}
\title{OrchLang: Auditing Resource Contracts for Bounded LLM Workflows\tnoteref{draft}}
\tnotetext[draft]{Author-review draft, 24 October 2026. Public archival release and final author declarations are pending. Not submitted.}
\author[ci]{Nambi Rajan M}
\ead{nambirajan.m2024@vitstudent.ac.in}
\author[analytics]{Ranjithkumar S\corref{cor}}
\ead{ranjithkumar.s@vit.ac.in}
\cortext[cor]{Corresponding author.}
\address[ci]{Department of Computational Intelligence, SCOPE, Vellore Institute of Technology, India}
\address[analytics]{Department of Analytics, SCOPE, Vellore Institute of Technology, India}
\begin{abstract}
OrchLang is an experimental compiler and audit artifact for bounded workflows using language models. It separates synthetic resource bounds, request equality and execution assumptions. An audit found an accepted workflow whose secret branches consumed 2 and 987 mock units. The repaired implementation distinguishes bindings from semantic contracts and retains request content. Separate readers reconstruct closed and response-dependent fragments for Lean checking. A text-valued core covers downstream requests under shared public inputs, secret Boolean control and stateless kernels. Reproduction includes 576 closed-request programs, 256 pipelines with 3,072 execution contexts, tokenizer probes and 11 documented external cost projections. These support scoped engineering claims, not provider privacy or a new general cost logic.
\end{abstract}
\begin{keyword}
resource analysis \sep information flow \sep workflow languages \sep certificate checking \sep reproducibility
\end{keyword}
\end{frontmatter}
\clearpage

\section*{Metadata}
\begin{table}[!ht]
\small
% Description labels follow the official template; widths fit the class margins.
\begin{tabular}{|l|p{.38\linewidth}|p{.47\linewidth}|}
\hline
\textbf{Nr.} & \textbf{Code metadata description} & \textbf{Please fill in this column} \\
\hline
C1 & Current code version & Publication audit candidate; current core \texttt{00420ef}; earlier results \texttt{6f325be} and \texttt{7a0095d}. Full hashes in manifests. \\
\hline
C2 & Permanent link to code/repository used for this code version & \href{https://github.com/guyoverclocked/orchlang-compiler-design-project}{Project repository}. Tested revision supplied locally; archival release pending. \\
\hline
C3 & Permanent link to Reproducible Capsule & Not yet deposited. Local source bundles and raw records supplied with draft. \\
\hline
C4 & Legal Code License & MIT (software); LPPL for Elsevier template derivative. Third-party notices included. \\
\hline
C5 & Code versioning system used & Git \\
\hline
C6 & Software code languages, tools, and services used & C++17, Python, Lean 4.19.0; optional Tokenizers 0.21.4. No inference service. \\
\hline
C7 & Compilation requirements, operating environments and dependencies & CPU, C++17 compiler, Python 3.10+. Tested: GCC 16.1, Python 3.14.3, Windows 11. Tokenizer probe: Python 3.12.14. \\
\hline
C8 & If available, link to developer documentation/manual & Repository README and included \texttt{docs/REPRODUCIBILITY.md}, \texttt{SEMANTICS.md}, \texttt{CERTIFICATES.md}. \\
\hline
C9 & Support email for questions & \texttt{ranjithkumar.s@vit.ac.in} \\
\hline
\end{tabular}
\caption{Code metadata. Public deposition remains a release requirement.}
\end{table}

\section{Motivation and significance}
A resource bound limits consumption; it need not make consumption independent
of a secret. A proof about a request kernel also requires an adapter that
serializes that request and accounts for billed failures. OrchLang provides a
small executable setting for inspecting and falsifying such assumptions.

The original checker at commit \texttt{a49fcdb} accepted a secret conditional
whose branches declared models named \texttt{m} with output caps 1 and 1000.
It equated their signatures by name. With seed 1 and identical public pins,
the branches consumed 2 and 987 synthetic units. All 95 original tests passed.
This refutes the original acceptance guarantee, although the unary bound of
1001 covers both observations. The artifact retains the witness and raw logs.

Resource-aware noninterference already combines information flow and cost
analysis~\cite{ngo}. RelCost reasons about relative costs~\cite{relcost}, and
probabilistic AARA derives expected-cost bounds~\cite{wang}. Randomness does
not preclude numeric specifications. eRHL supplies a richer relational
probabilistic logic~\cite{erhl}. Hard-KLAIM proves typed trace safety under
interface assumptions~\cite{hard}; PC2 checks constraints before operations
and identifies conditions for real-world correspondence~\cite{pc2}. Loom's
author-hosted report addresses structural LLM workflow cost bounds~\cite{loom};
its peer-review status and artifact results are not established here.
OrchLang claims no expressiveness or priority separation. Its software
contribution is an inspectable specialization with counterexamples, repaired
identity tracking, a separate certificate path and reproducible validation.

\section{Software description}
\subsection{Architecture and functionality}
The C++ frontend supports bounded inputs, prompts, lexical models, conditionals,
bounded retries, outputs and tool sinks. Analyses track confidentiality,
integrity and synthetic resource bounds. Resolved binding identities prevent
shadowing errors. Canonical abstract contracts admit equivalent aliases while
distinguishing caps and model revisions. Relational signatures retain template
text, ordered typed arguments, literal content and resolved public guards.
A secret branch is conservatively rejected when its request relation cannot
be certified.

Commands \texttt{check}, \texttt{ir-json}, \texttt{certify} and \texttt{run}
produce diagnostics, structural IR, an analysis report and length-only mock
execution. Despite its historical name, \texttt{certify} emits an explicitly
labelled report. It is not independently checked proof evidence. Numeric
parsing rejects malformed pins and overflow; runtime pins must respect declared
bounds. Arithmetic uses safe division and saturation, with saturated static
totals rejected.

\begin{figure}[!ht]
\centering
\fbox{\begin{minipage}{.92\linewidth}
Source $\longrightarrow$ C++ analysis/report and length-only mock.\\[3pt]
Same source, restricted fragments $\longrightarrow$ independent Python readers
and digest binding $\longrightarrow$ JSON cores and untrusted plans
$\longrightarrow$ Lean decoder/checker.\\[3pt]
Hand-authored text core $\longrightarrow$ finite reference interpreter
$\longleftrightarrow$ separate finite-table replay executor.\\[3pt]
Trusted boundaries: source reader/serialization, JSON library, primitive
contracts, tokenizer library and execution platform.
\end{minipage}}
\caption{Separate evidence paths. C++ acceptance is not a premise of the Lean
theorem; no verified translation covers all source features.}
\end{figure}

Separate certificate paths independently reparse source, reconstruct request
expressions and check source digest/core agreement before invoking Lean on an
untrusted plan. The original path admits closed requests. The new dataflow path
also admits shared public text/Boolean inputs and response bindings used by
later requests. Source-local names are aligned to lexical slots. Secret text,
source retries, tools and reclassification fail closed. Dataflow outputs must
be root literal statements. The Python readers and slot allocation are trusted,
not mechanically verified. Neither path checks the C++ mock budget as a
tokenizer budget.

\subsection{Formal guarantees and assumptions}
The unary Lean core bounds scalar cost for calls, sequence, branch maxima and
up-to-$n$ repetition, assuming primitive call bounds. The closed relational core
uses a shared stateless kernel indexed by contract and complete request. The
dataflow core additionally evaluates literal, slot and append expressions in a
text store. A call binds actual response text; subsequent requests use that
store. Its kernel $K(c,r)$ returns text, status and weight. Events record contract,
UTF-8 input/output byte counts and failure status. Sequence composes stores,
concatenates traces and multiplies weights. Lexical bounded retry restarts from
the outer store, retains failures and stops after an attempt with no failed
call. Errors are status values; the body completes before retry decides. This
is not an SDK exception semantics or the C++ mock's independent failure sample.

Normalization removes a secret branch only when its arms have exactly the
same plan. Lean's \texttt{normalize\_correct} proves
\[
 N(P)=\operatorname{some}(Q)\ \Longrightarrow\
 \mathcal{D}_{K,l,h}(P)=\mathcal{D}^{\mathrm{plan}}_{K,l}(Q)
 \quad\text{for every Boolean store }h.
\]
Here $l$ is a shared public Boolean store. In the dataflow core the equality is
between complete text-state transformers, for every shared initial text store.
Thus secret Boolean stores yield equal weighted observations even when later
requests depend on earlier response contents. A deterministic observer preserves
equality. The result is algebraic; probability interpretation additionally
requires normalized nonnegative finite kernels. Runtime checks validate the
supplied tables, not a provider's compliance.

The executable decoder/checker has the theorem
\path{decoded_certificate_sound} for successfully decoded core values.
Lean 4.19.0 checks all three files without admitted lemmas or custom axioms;
reported axioms are \texttt{propext} and \texttt{Quot.sound}. Complete source
correspondence, secret text inputs and stateful providers remain outside the proof.
Hash binding detects source changes; it authenticates neither authorship nor
the reader's correctness.

\section{Illustrative examples}
Requests \texttt{yes} and \texttt{no!} have the same byte length. A valid finite
kernel may return one byte for the former and 100 for the latter. Input-length
equality cannot justify equal output distributions; a common random-length
tape would alter a marginal. The repaired checker retains content and rejects
this example. Equivalent model aliases with equal constructed requests are
accepted in \texttt{examples/valid/closed\_certificate.orch}. The certificate
manual provides executable issuance/checking commands.

A downstream witness returns \texttt{aa} or \texttt{bb} with equal probability.
Their lengths agree, but requests \texttt{Next:aa} and \texttt{Next:bb} produce
one and three output bytes. Replacing the response-dependent request with the
constant \texttt{Next:aa} changes its distribution. The dataflow checker rejects
that substitution and admits two matched pipelines with renamed local bindings.
Its source example is in \path{examples/certificates/dataflow_certificate.orch}.

Ordered observations contain billed events. Aggregate observations sum usage,
calls and failures by contract and unit. Reversing different calls can preserve
the aggregate but change the ordered trace. Timing, provider-visible contents
and tool effects are excluded. Equal bills do not authorize secret disclosure.
For a call succeeding with probability $1/2$, two-attempt retry yields success
with weight $1/2$, failure then success with weight $1/4$, and two failures with
weight $1/4$. Failed attempts remain billed. Lean guards and replay observations
check this example.

\section{Impact and evaluation}
OrchLang supports reproducible investigation of contract omissions. Evidence
is engineering validation on a developer machine; independent human reproduction,
production use and external adoption are not established. Clean local snapshots,
source hashes, commands, raw observations and derived tables accompany the
draft. CPU checks use GCC 16.1, C++17, warnings enabled and optimization level 2.
No GPU, model inference or paid API was used. All counts below are generated
from retained records by \path{scripts/report_publication.py} and
\path{scripts/report_dataflow.py}.

The current suite passes \evDataflowCppTests{} C++ and \evDataflowPythonTests{} Python tests.
Four compiled semantic mutants are detected: name conflation, request-content
erasure, missing retry multiplication and unbilled failure. A structural
baseline using maximum, sum and retry multiplication agrees on \evCostWorkflows{} existing
cost workflows and \evCostRuns{} mock runs. It shares frontend primitive facts,
which it cannot independently validate. Historical baselines omitting retries
are retained only as broken negative controls.

Exhaustive development enumeration combines 24 arm variants into
\evGenerated{} programs, with four call-count structures. The compiler accepts
\evAccepted{}, with \evAcceptedFiniteCounterexamples{} counterexamples under
the supplied finite kernels. Decisions match conventional trace identity
exactly. Sixteen rejected programs are equivalent under those particular
kernels; they are not universally safe. No precision advantage is established.
The separate source reader and Lean path perform \evLeanChecks{} checks, issue
\evCertificates{} certificates and reject \evSourceTampers{} source changes,
24 altered plans and four malformed batches. Reference and replay observations
agree on \evReplayCases{} cases comprising \evReplayPaths{} weighted paths.
Enumeration is not a universal soundness proof.

The new development grammar covers \evDataflowSourcePrograms{} source programs
with one two-call-per-arm pipeline shape and a public optional tail. Three text
inputs, two public flags and two secret values give
\evDataflowExecutionContexts{} contexts. The independently constructed expected
cores match the source reader throughout. Python and Lean normalization agree;
Lean execution matches reference event distributions and replay request/event
distributions in every context. The normalizer admits \evDataflowAccepted{}
programs, with \evDataflowAcceptedFiniteCounterexamples{} supplied-kernel
counterexamples. Eight additional core cases, comprising
\evDataflowBoundaryPaths{} paths, cover UTF-8, nested/zero retry and the downstream
witness. Tests reject \evDataflowSourceTampersRejected{} changed sources,
\evDataflowAlteredPlansRejected{} wrong plans and
\evDataflowInvalidEvaluatorInputsRejected{} invalid evaluator inputs.

The C++ checker accepts \evDataflowCppAccepted{} of those source cases and emits
E230 in \evDataflowCppSecretArgumentCases{}: its PC label marks branch-local responses
secret. The dataflow resource judgment does not override that confidentiality
policy; the counts are not a precision ranking. Initial campaign attempts failed
on a transport-size guard and unused prompt parameters. Batching and zero-argument
fixed prompts repaired them without changing expected requests/cores or kernels.
Failed records and amendments are retained. This is development validation,
not held-out evaluation or a new general logic.

Tokenizer development probes use pinned GPT-2 BPE and BERT WordPiece files with
Tokenizers 0.21.4. Of \evTokenizerEncodings{} fixed encodings,
\evUnderestimates{} exceed the estimate $\lceil b/4\rceil$ for $b$ UTF-8 bytes;
\evNonadditive{} of \evConcatenations{} concatenations differ from summed
component counts. All token IDs and special-token options are retained. These
cases refute a general tokenizer-bound interpretation, not estimate a failure
rate. The C++ unit is explicitly \texttt{mock-length-v1}.

An external census considers 15 Prompt Flow examples and six LangGraph tutorial
patterns at pinned revisions~\cite{promptflow,langgraph}. Eleven authored cost
projections are included; ten exclusions have recorded reasons. Three basic
configurations share a call structure. One helper named \texttt{llm\_result}
returns a fixed string upstream and contributes zero model calls. Projections
introduce bounded inputs/caps, abstract guards, summarize prompts and may
serialize concurrency or omit aggregation. They preserve only documented cost
shape: zero original applications or faithful end-to-end translations were
executed. The corpus has no secret branches and supplies no relational
precision evidence.

\begin{table}[!ht]
\centering\small
\begin{tabular}{lrrrrr}
\hline
Projection & Calls & Input & Output & Total & Runs \\
\hline
pf-basic & 1 & 260 & 120 & 380 & 20 \\
pf-basic-with-builtin-llm & 1 & 260 & 120 & 380 & 20 \\
pf-basic-with-connection & 1 & 260 & 120 & 380 & 20 \\
pf-customer-intent & 1 & 392 & 256 & 648 & 20 \\
pf-ner & 1 & 533 & 64 & 597 & 20 \\
pf-switch & 1 & 262 & 128 & 390 & 20 \\
pf-if-else-stubs & 0 & 0 & 0 & 0 & 20 \\
pf-web-classification & 2 & 2203 & 256 & 2459 & 20 \\
lg-prompt-chaining & 3 & 402 & 192 & 594 & 20 \\
lg-parallelization & 3 & 786 & 192 & 978 & 20 \\
lg-routing & 2 & 522 & 128 & 650 & 20 \\
\hline
\end{tabular}

\caption{External cost projections in synthetic mock units. Compiler and
structural bounds agree throughout. Changes/exclusions are in the manifest.}
\label{tab:external}
\end{table}

The protocol fixes seeds 1--20 and bound/call-count endpoints. Its initial run
failed on a reserved identifier in one projection; this was corrected and an
amendment frozen before the complete rerun. Partial-record persistence was also
repaired. The failure record is retained; this is not untouched preregistration.
The amended run completes \evExternalRuns{} mock executions with
\evExternalBoundViolations{} componentwise violations and
\evExternalBaselineDisagreements{} baseline disagreements
(Table~\ref{tab:external}). Repeated seeds are not independent application
samples. No significance, throughput or provider-conformance claim follows.

\section{Conclusions}
Confusing names or lengths with semantic requests can produce false resource
certifications despite passing tests. OrchLang makes this failure reproducible
and supplies repaired analysis, separate source-bound plans, scoped Lean results
and replay checks. Evidence supports synthetic bounds and text-dataflow equality
under explicit contracts. General source verification, stateful providers and
inference adapters remain open. Measuring annotation effort and retained
functional behavior in independently adapted workflows is a next evaluation;
the current projections do not answer that question.

\section*{Data, software and research-assistance statements}
Source, raw records, protocols, claim ledger, reproduction instructions and
tested-source bundles accompany the draft. Private library PDFs, credentials
and model weights are excluded. A public versioned archive remains required
before submission. The author reports no funding. Competing-interest
declarations and both authors' final approval await confirmation.

OpenAI Codex assisted code, proofs, literature retrieval, scripts and writing.
Its research-method role included constructing
and testing the artifact. Both authors must verify the code, citations, proofs,
results and text before submission; completed human review is not asserted.

\begin{thebibliography}{10}
\bibitem{ngo} V. C. Ngo, M. Dehesa-Azuara, M. Fredrikson, J. Hoffmann,
Verifying and Synthesizing Constant-Resource Implementations with Types,
IEEE Symposium on Security and Privacy (2017).
\url{https://www.ieee-security.org/TC/SP2017/papers/408.pdf}.
\bibitem{relcost} E. \c{C}i\c{c}ek, G. Barthe, M. Gaboardi, D. Garg, J. Hoffmann,
Relational Cost Analysis, POPL (2017). \url{https://doi.org/10.1145/3009837.3009858}.
\bibitem{wang} D. Wang, D. M. Kahn, J. Hoffmann, Raising Expectations:
Automating Expected Cost Analysis with Types, Proc. ACM Program. Lang. 4,
ICFP, Article 110 (2020). \url{https://doi.org/10.1145/3408992}.
\bibitem{erhl} M. Avanzini, G. Barthe, D. Davoli, B. Gr\'egoire,
A quantitative probabilistic relational Hoare logic, arXiv:2407.17127v2 (2025).
\url{https://arxiv.org/abs/2407.17127v2}.
\bibitem{hard} M. Aldinucci, A. Bracciali, I. Colonnelli, D. Medi\'c,
A. Mulone, L. Roversi, M. E. Santimaria, Hard-KLAIM: Hardening MCP-Enabled
Agentic Workflows with a Locality-Aware Type-and-Effect System for Trace Safety,
Journal of Logical and Algebraic Methods in Programming, pre-proof (2026).
\url{https://doi.org/10.1016/j.jlamp.2026.101191}.
\bibitem{pc2} C. Ciliberto, A. Rudi, E. Tassi, Proof-Carrying Contracts:
Commitments AI Agents Cannot Break, arXiv:2610.28484v1 (2026).
\url{https://arxiv.org/abs/2610.28484v1}.
\bibitem{loom} K. Dickerson, Statically Bounding Resource Accumulation in LLM
Agent Workflows: The Cost Case, author-hosted report, 18 May 2026.
\url{https://loom.technology/research/bounded-llm-cost/}.
\bibitem{promptflow} Microsoft, Prompt Flow, standard examples, revision
\nolinkurl{33a68c9d3cf349d7fdfad1f41969af27543c0406}, accessed 23 October 2026.
\url{https://github.com/microsoft/promptflow}.
\bibitem{langgraph} LangChain, LangGraph workflows and agents documentation,
revision \nolinkurl{ecda7f522a6ffd5faaabb843b581fe57cd5eb782}, accessed 23 October 2026.
\url{https://github.com/langchain-ai/docs}.
\end{thebibliography}
\end{document}
